\subsection{REPcc Operation Examples}
The following examples show how the observed value, REPFLAGS, counter update, and terminating iteration interact
for common operation classes. In each case, the counter is tested before the body executes. The body then
computes an observed value, derives REPFLAGS from that value, and updates the counter only when the selected REP
condition remains true.

\paragraph{Counted Move.}
Plain REP is the REPT alias, so the condition is always true. This form copies exactly the requested count, and
the counter reaches zero when the transfer finishes.

\begin{quote}
\begingroup\small\ttfamily
\begin{tabular}{@{}l@{}}
; R0 = src, R1 = dst, R2 = count\\
REP R2, MOV.B [R0++], [R1++]\\
\end{tabular}
\endgroup
\end{quote}

For MOV, the observed value is the source value. REPT does not use it to stop; the body commits and the counter
moves toward zero after every completed iteration.

\paragraph{Scan Until Zero.}
MOV can also be used with REPNE to scan a stream while the loaded value remains nonzero.

\begin{quote}
\begingroup\small\ttfamily
\begin{tabular}{@{}l@{}}
; R0 = base, R2 = count\\
REPNE R2, MOV.L [R0++], R3\\
\end{tabular}
\endgroup
\end{quote}

The MOV body loads the candidate element into R3 and uses that loaded value as the observed value. If the value
is nonzero, REPNE is true and R2 is decremented. If the value is zero, the load and address update still commit,
but R2 is not decremented.

\paragraph{Compare While Equal.}
CMP uses the subtraction result as the observed value. REPEQ continues while the comparison is equal and stops
on the first non-equal comparison or when the counter reaches zero.

\begin{quote}
\begingroup\small\ttfamily
\begin{tabular}{@{}l@{}}
; R0 = a, R1 = b, R2 = count\\
REPEQ R2, CMP.B [R0++], [R1++]\\
\end{tabular}
\endgroup
\end{quote}

When a mismatch is found, the terminating CMP commits its FLAGS result and the counter is not updated for that
terminating iteration. Software can use the final FLAGS state to derive the ordered result.

\paragraph{Test While Nonzero.}
TEST derives REPFLAGS from the logical test result. A memory scan can therefore continue while a masked value
remains nonzero, or stop when a selected bit range clears.

\begin{quote}
\begingroup\small\ttfamily
\begin{tabular}{@{}l@{}}
; R0 = base, R2 = count, R3 = mask\\
REPNE R2, TEST.L [R0++], R3\\
\end{tabular}
\endgroup
\end{quote}

The memory operand and mask are tested each iteration. A nonzero tested result updates the counter and
continues; a zero tested result commits the TEST FLAGS and leaves the counter at the terminating element.

\paragraph{Bulk Extension.}
The EXTS and EXTZ memory forms use the widened result as the observed value. An unconditional REP form expresses
a compact bulk conversion, while a conditional form can stop when a converted value matches the selected
condition.

\begin{quote}
\begingroup\small\ttfamily
\begin{tabular}{@{}l@{}}
; R0 = byte source, R1 = long destination, R2 = count\\
REP R2, EXTZL.B [R0++], [R1++]\\
\end{tabular}
\endgroup
\end{quote}

The source byte is read, zero-extended to a long value, and written to the destination element. Because REP is
REPT, every completed conversion updates the counter.
